% =====================================================================
%  Chapter 16 — Generalised Projections and the \idxw{Uniqueness} Theorems
% =====================================================================

\chapter{Generalised Projections and the Uniqueness Theorems}
\label{ch:projections}

Every shape measurement in this book is a projection. A photograph of a
rotating body is the body's silhouette against the sky, taken along a
direction that changes as the body turns. A radar image is a projection\idx{Projection} of a
different kind, reconstructed from the delay and the Doppler shift of an
echo. An interferometric image is a projection\idx{Projection} reconstructed from the
visibility and the phase of a fringe pattern.

This chapter asks the question that underlies all of them: what can be
determined from a projection\idx{Projection}, and how many of them are needed?

\section{What a single projection gives}

A projection along a fixed direction of a convex body\idx{Convex body} gives the outline of
its shadow: a convex region whose area is the body's projected area and
whose shape is the body's projected shape. From that one outline, three
things follow immediately --- the projected area, the projected extent in
two orthogonal directions, and the rotation phase at which the projection
was taken.

Three things do not follow, and the list is worth keeping in view because it
is the whole of Chapter 8's argument:

The volume does not follow. A sphere and a flat disk of the same projected
area have the same silhouette and volumes differing by a factor of three.
The projected area constrains the volume only to within the range that the
convex hull of the silhouette allows.

The three-dimensional shape does not follow. A body of revolution\idx{Body of revolution} and an
ellipsoid with the same silhouette are indistinguishable in a single
projection, and so are a thousand non-convex bodies whose shadow agrees.

The orientation of the body in space does not follow, beyond the two
projected axes. The direction of the spin pole is what fixes it, and a
single projection does not fix the pole.

\section{Why the classical uniqueness results matter}

The theory of reconstructing a body from its projections is old and precise,
and three of its results are directly relevant.

A convex body\idx{Convex body} is uniquely determined by the set of all its orthogonal
projections. This is the strongest statement and the least useful, since it
requires infinitely many.

A body of revolution\idx{Body of revolution} is determined by a single projection along its axis of
symmetry. This is the strongest useful statement, and it is why the
rotation state of an asteroid is such valuable information: a body that is
close to a body of revolution\idx{Body of revolution} can be characterised from one projection, and
the residual error is a measure of how far from a body of revolution it is.

And a general body requires a family of projections. For a convex body\idx{Convex body}, the
projections onto a small number of suitably chosen directions --- in the
limit, three generic directions for a polytope --- can determine it, but
``suitably chosen'' is a real constraint: the directions must not share a
degenerate relationship, and a single direction never suffices.

\section{What the rotation supplies}

This is where the astrophysical case differs from the geometric one, and the
difference is favourable in a way that is not obvious before it is stated.

A body on the ground is observed from essentially one direction, and that
direction is fixed by where the observer is. As the body rotates, the
\emph{body's} orientation changes but the \emph{projection direction} does
not. What changes is the shape of the projected outline, because the body
presents different parts of itself to the observer.

That is the whole content of the lightcurve. The brightness of a rotating
body varies because the projected outline and the projected brightness
distribution vary. The variation has a period, and the period is the
rotation period.

So the body does supply a family of projections --- the set of all its
projections over a full rotation. For a body whose projection changes
sufficiently with phase, the set of projections over one rotation carries
the information of a band of directions, and it is that band, rather than a
single outline, that the inversion uses.

\begin{deeper}{Why the inversion is nevertheless not unique}
The band is not the whole sphere of directions. The projections obtained
over one rotation all share the same viewing direction and differ only in
the body's orientation about one axis, so the reconstruction is a body of
revolution about that axis --- up to the non-convexity, which the
projections at different phases can partly expose. The practical result is
the familiar one: a single lightcurve admits a family of shapes and a
family of pole directions, and the standard inversion returns the convex
hull of that family. Additional lightcurves at different apparitions, or a
single resolved image that fixes the geometry directly, are what break the
degeneracy.
\end{deeper}

\section{What this implies for the measurement hierarchy}

The hierarchy of methods in the field follows from this, and it explains why
the cheapest methods produce the least certain results.

\begin{table}[!htbp]
\centering
\small
\caption{What each measurement family determines, and the assumption each
one carries.}
\label{tab:hierarchy}
\begin{tabularx}{\textwidth}{@{}P{1.24in}P{1.32in}Y@{}}
\toprule
\rowcolor{panel}
\textbf{Method} & \textbf{Determines} & \textbf{Carries} \\
\midrule
Lightcurve & period, pole, low-order shape & convexity; the band of
projections is not the whole sphere \\
\rowcolor{panel}
Disk image & the projected outline at one epoch & the shape of that
outline; the volume still needs a model \\
Radar image & the three-dimensional shape directly & the geometry of the
measurement \\
\rowcolor{panel}
Interferometric image & the projected outline at high angular resolution &
the visibility model, and the matching of the baseline to the target \\
\bottomrule
\end{tabularx}
\end{table}

The last row is where the book's argument goes next. An interferometric
image is, in principle, the best of the four, because the angular resolution
it delivers is a function of the baseline rather than of the target's
distance, and because the resolution can be set by choosing the baseline
rather than by waiting for a close approach.

The catch is that an interferometer does not produce a picture by itself. It
produces a visibility and a phase, and the picture is reconstructed from
those, and the reconstruction has an assumption attached: a model of what
the body's brightness distribution looks like on the sky. That model is the
subject of Chapter 20.

\begin{keytakeaways}
\item A single projection determines the projected area and the projected
extent, and does not determine the volume, the three-dimensional shape or
the orientation.
\item The rotation supplies a band of projections rather than a single one,
which is why a lightcurve carries more information than a single image ---
and why it still returns a convex hull.
\item The classical uniqueness\idx{Uniqueness} results set the target: a body of revolution
is determined by one axial projection, and a general convex body needs a
family.
\item An interferometer's resolution is set by the baseline rather than by
the distance, which is why it is the only method that can image a small
object without waiting for a close approach.
\end{keytakeaways}
